\scanpage{5}
\mspage{3}
ou si on préfère
\[
H^1(X_{\overline K},\mathbb Z_\ell)\simeq T_\ell(\underline{\operatorname{Pic}}_{X/K}(\overline K))\underset{\mathbb Z_\ell}{\otimes}\mathbb Z_\ell(-1).
\]

0.6. Lorsque \(A\) est un schéma abélien sur \(K\), nous verrons que la connaissance du \(\pi\)-module \(T_\ell(A(\overline K))\) donne des informations \unsure{assez} précises sur la structure du \uline{modèle de Néron} de \(A\) [\quad], -- et inversement la \inserted{considération du mod. de Néron permet d'obtenir des renseignements complémentaires intéressants sur \(T_\ell(A(\overline K))\), -- et partant sur tous les \(H^1(X_{\overline K})\unsure{\bar s|;|f}\), comme \(\pi\)-modules.}
% slanted: continuation of the insertion in the left margin

D'autre part, l'étude des \(\pi\)-modules plus généraux \(H^i(X_{\overline K},\mathbb Z_\ell)\) s'introduit quand on \deleted{\unsure{désigne de}} dispose d'un \deleted{modèle} \inserted{schéma} \(X\) \deleted{\unsure{de}} sur \(S\) de fibre générique \(X_K=X_\eta\), et qu'on désire \deleted{\unsure{étudier}} les relations entre \inserted{\(H^i(X)\),} \(H^i(X_{\unsure{s|\bar s}})\) et \(H^i(X_{\bar\eta})\) ; (p. ex. dans certains cas on pourra montrer \(H^i(X_s)\simeq H^i(X_{\bar\eta})^\pi\)). Cette \deleted{Deligne, utilisant l'étude} étude locale s'avère comme un outil indispensable pour l'étude \inserted{cohomologique} \uline{globale} d'un morphisme
\[
f:X\mathrel{\unsure{-|\longrightarrow}}S
\]
avec \(S\) régulier de dim 1. \deleted{\unsure{Prenant les|Poussant les|Prenant le|Pour}} Dans le cas où \(S=\mathbb P^1_k\), et où \(\deleted{X}f\) ``est'' \deleted{\unsure{déduit|décrit}} un pinceau \inserted{\unsure{linéaires|linéaire}} de sections hyperplanes d'une variété projective lisse \(X_0\) définie sur \(k\), on est dans la situation typique de la \uline{théorie de Lefschetz} [\quad]. S'inspirant d'arguments de Lefschetz, Deligne montrera plus tard que si \(X_0\) est une surface de degré \(\geq4\)
% archive sheet number: 4
